\begin{definition}
Definicja iloczynu kartezjańskiego dwóch zbiorów: $\\$
Niech $A$ i $B$ będą dowolnymi zbiorami. Wówczas $A\times B=\{(x,y): x\in A\wedge y\in B\}$. $\\$
Zatem iloczyn kartezjański zbiorów $A$ i $B$ to zbiór par uporządkowanych $(x,y)$, gdzie $x\in A$ i $y\in B$. $\\$
Warunek przynależności pary elementów $(x,y)$ do iloczynu kartezjańskiego zbiorów $A$ i $B$: $(x,y)\in A\times B\Leftrightarrow x\in A\wedge y\in B$. $\\$
Iloczyn kartezjański zbiorów jest relacją. Jeżeli $R\subseteq X\times Y$ to $R$ jest relacją dwuargumentową.
Zatem każdy zbiór $R\subseteq X\times Y$ jest relacją. Stąd relacja jest podzbiorem zbioru wszystkich możliwych
par uporządkowanych. Jeżeli $R\subseteq X\times Y$ to mówimy, że $R$ jest relacją pomiędzy $X$ i $Y$. Dodajmy, że $R=X\times Y$ jest
relacją pełną pomiędzy $X$ i $Y$. Zauważmy też, że zbiór pusty jest relacją. $\forall(x,y)[(x,y)\in\emptyset\Rightarrow (x,y)\in X\times Y]\Leftrightarrow \emptyset\subseteq X\times Y$.
Jest to tzw. implikacja próżniowa. Jej poprzednik jest zawsze fałszywy. Zatem implikacja próżniowa jest zawsze prawdziwa, ponieważ z definicji implikacja, której poprzednik jest fałszywy jest zawsze prawdziwa. 
Niezależnie od jej następnika. Zauważmy, że w powyższej implikacji nawet jeśli $X$ i $Y$ są puste to ta implikacja jest prawdziwa. Bo z fałszu wynika prawda albo fałsz.
\end{definition}
\begin{remark}
Zaprzeczenie definicji przynależności pary elementów $(x,y)$ do iloczynu kartezjańskiego zbiorów $A$ i $B$:
$\neg(x,y)\in A\times B\Leftrightarrow\neg(x\in A\wedge y\in B)\Leftrightarrow\neg(x\in A)\vee\neg(y\in B)\Leftrightarrow x\notin A\vee y\notin B$.
\end{remark}
\begin{property}
Dla niepustych zbiorów $X,Y,A$ i $B$ zachodzi implikacja $A\times B\subseteq X\times Y\Rightarrow A\subseteq X\wedge B\subseteq Y$.
\end{property}
\begin{proof}[\textbf{Dowód}]
Ustalmy dowolne, niepuste zbiory $X,Y,A$ i $B$ oraz dowolne elementy $x$ i $y$. Niech $x\in A$, $y\in B$ i $A\times B\subseteq X\times Y$. Wówczas z definicji $(x,y)\in A\times B$. Ale z założenia mamy $(x,y)\in X\times Y$. Stąd z definicji $x\in X$ i $y\in Y$. Ale wobec tego, że $x$ i $y$ zostały wybrane jako dowolne elementy to prawdziwe są implikacje $x\in A\Rightarrow x\in X$ i $y\in B\Rightarrow y\in Y$, a wówczas z definicji zawierania się zbiorów otrzymujemy, że $A\subseteq X$ i $B\subseteq Y$. Zatem zadana implikacja zachodzi dla niepustych zbiorów.
\end{proof}
\begin{property}
Dla dowolnych zbiorów $X,Y,A, B$ zachodzi implikacja
$A\subseteq X\wedge B\subseteq Y\Rightarrow A\times B\subseteq X\times Y$.
\end{property}
\begin{proof}[\textbf{Dowód}]
Ustalmy dowolne zbiory $X,Y,A\subseteq X, B\subseteq Y$ oraz dowolne elementy $x$ i $y$. Niech
$x\in A$ i $y\in B$. Wówczas z definicji $(x,y)\in A\times B$. Ale skoro $A\subseteq X$ to $x\in X$ i $B\subseteq Y$ to $y\in Y$. Stąd $(x,y)\in X\times Y$. Ponieważ $x$ i $y$ zostały wybrane dowolnie to zadana implikacja zachodzi w ogólności.
\end{proof}
\begin{property}
Dla dowolnych zbiorów $X,Y, A$ i $B$ nie zachodzi implikacja $A\times B\subseteq X\times Y\Rightarrow A\subseteq X\wedge B\subseteq Y$.
\end{property}
\begin{proof}[\textbf{Dowód}]
Ustalmy dowolne zbiory $A=\emptyset$, $B=\{y_{1},y_{2}\}$ oraz $X=\{x_{3},x_{4},x_{5}\}$ i $Y=\{y_{3},y_{4},y_{5}\}$, gdzie $y_{1},y_{2}\ne y_{3},y_{4},y_{5}$. Wówczas warunki $A\times B\subseteq X\times Y$ i $A\subseteq X$ są spełnione ale $B\not\subseteq Y$. Zatem zadana implikacja nie jest w ogólności prawdziwa.
\end{proof}
\begin{property}
Dla niepustych zbiorów $A,B,C$ i $D$ zachodzi równoważność:
$A\times B=C\times D\Leftrightarrow A=C\wedge B=D$.
\end{property}
\begin{proof}[\textbf{Dowód}]
Ustalmy dowolne elementy $x$ i $y$ oraz dowolne niepuste zbiory $A,B,C$ i $D$. Niech $x\in A$ i $y\in B$. Wówczas z definicji iloczynu kartezjańskiego zbiorów wynika, że $(x,y)\in A\times B$. Ale z równości $A\times B=C\times D$ otrzymujemy, że $(x,y)\in C\times D$. Stąd $x\in C$ i $y\in D$. Skoro dla każdego $x\in A$ jest $x\in C$ i dla każdego $y\in B$ jest $y\in D$ to $A\subseteq C$ i $B\subseteq D$. Zatem prawdziwa jest implikacja $A\times B\subseteq C\times D\Rightarrow A\subseteq C\wedge B\subseteq D$. $\\$
Ustalmy dowolne elementy $x$ i $y$ oraz dowolne niepuste zbiory $A,B,C$ i $D$. Niech $x\in C$ i $y\in D$. Wówczas z definicji $(x,y)\in C\times D$. Ale z równości $A\times B=C\times D$ mamy, że $(x,y)\in A\times B$. Stąd $x\in A$ i $y\in B$. Ponieważ $x$ i $y$ zostały wybrane dowolnie to dostajemy, że jeżeli $x\in C$ to $x\in A$ i jeżeli $y\in D$ to $y\in B$ czyli $C\subseteq A$ i $D\subseteq B$. Zatem prawdziwa jest implikacja $C\times D\subseteq A\times B\Rightarrow C\subseteq A\wedge D\subseteq B$. Wówczas zachodzi $A\times B=C\times D\Rightarrow A=C\wedge B=D$. $\\$
Ustalmy dowolne elementy $x$ i $y$ oraz dowolne niepuste zbiory $A,B,C$ i $D$. Niech $x\in A$ i $y\in B$. Wówczas z definicji $(x,y)\in A\times B$. Ponieważ $A=C$ i $B=D$ to $x\in C$ i $y\in D$. Zatem $(x,y)\in C\times D$. Wówczas $A\times B\subseteq C\times D$. Zatem prawdziwa jest implikacja $A=C\wedge B=D\Rightarrow A\times B\subseteq C\times D$. $\\$
Ustalmy dowolne elementy $x$ i $y$ oraz dowolne niepuste zbiory $A,B,C$ i $D$. Niech $x\in C$ i $y\in D$. Wówczas z definicji $(x,y)\in C\times D$. Ale $C=A$ i $D=B$. Stąd $x\in A$ i $y\in B$. Zatem $(x,y)\in A\times B$. Ponieważ jeżeli $(x,y)\in C\times D$ to $(x,y)\in A\times B$ to dostajemy $C\times D\subseteq A\times B$. Wówczas prawdziwa jest implikacja $A=C\wedge B=D\Rightarrow C\times D\subseteq A\times B$. $\\$
Wobec prawdziwości wszystkich czterech implikacji otrzymujemy, że zadana równość zachodzi dla niepustych zbiorów. $\\$
Zauważmy też, że zadana równoważność jest w szczególności prawdziwa, gdy wszystkie zbiory $A,B,C$ i $D$ są puste. W przypadku, gdy $A=\emptyset$ to $A\times B=\emptyset$. Wówczas $C\times D=\emptyset$. A zatem $C=\emptyset$ lub $D=\emptyset$. Wówczas mamy następujące możliwości:$\\$
$A=B=C=D=\emptyset$. Wtedy zadana równoważność jest prawdziwa. $\\$
$A=\emptyset$, $B,D\ne\emptyset$. Zatem musi być $C=\emptyset$. Wówczas $B=D$ lub $B\ne D$. Równość $A\times B=C\times D$ zachodzi bez względu czy $B=D$ czy $B\ne D$. Natomiast zadana równoważność zachodzi, gdy $B=D$ zaś nie zachodzi, gdy $B\ne D$. $\\$
$A=B=\emptyset$. Wówczas $C\times D=\emptyset$, a zatem $C=\emptyset$ lub $D=\emptyset$. Wtedy zachodzi równość $A\times B=C\times D$. Ale jeśli, któryś ze zbiorów $C$ lub $D$ jest niepusty to zadana równoważność nie jest spełniona.
Itd. $\\$
Powyższe rozważania prowadzą nas do kolejnej własności iloczynu kartezjańskiego zbiorów. Jej dowód poniżej.
\end{proof}
\begin{property}
Dla dowolnych zbiorów $A,B,C$ i $D$ zachodzi implikacja:
$A=\\=C\wedge B=D\Rightarrow A\times B=C\times D$.
\end{property}
\begin{proof}[\textbf{Dowód}]
Ustalmy dowolne elementy $x$, $y$ oraz dowolne zbiory $A,B,C$ i $D$. Niech $x\in A$ i $y\in B$. Wówczas z definicji wynika, że $(x,y)\in A\times B$. Ale z założenia $A=C$ i $B=D$. Stąd $(x,y)\in C\times D$ co wobec dowolności wyboru elementów $x$ i $y$ daje nam $A\times B\subseteq C\times D$. Wówczas prawdziwa jest implikacja $A=C\wedge B=D\Rightarrow A\times B\subseteq C\times D$. $\\$
Ustalmy dowolne elementy $x$, $y$ oraz dowolne zbiory $A,B,C$ i $D$. Niech $x\in C$ i $y\in D$. Wówczas z definicji wynika, że $(x,y)\in C\times D$. Ale z założenia $C=A$ i $D=B$. Zatem $(x,y)\in A\times B$. Ponieważ elementy $x$ i $y$ zostały wybrane dowolnie to otrzymujemy, że $C\times D\subseteq A\times B$. Stąd prawdziwa jest implikacja $A=C\wedge B=D\Rightarrow C\times D\subseteq A\times B$. Ponieważ zachodzą obie inkluzje to zadana implikacja jest prawdziwa w ogólności.
\end{proof}
\begin{property}
Niech $A,B,C,D$ będą zbiorami. W ogólności nie zachodzi implikacja $A\times B=C\times D\Rightarrow A=C\wedge B=D$.
\end{property}
\begin{proof}[\textbf{Dowód}]
Niech $A=\emptyset$, $B=\{x_{1},x_{2}\}$. Wówczas $A\times B=\emptyset$, a zatem $C\times D=\emptyset$. Niech $C=\{x_{3},x_{4}\}$. Przy czym $B\cap C=\emptyset$. Wówczas poprzednik zadanej implikacji jest prawdziwy. Ale $A\ne C$. Co wystarczy do fałszywości następnika zadanej implikacji. Stąd zadana implikacja nie zachodzi w ogólności.
\end{proof}
\begin{remark}
Patrząc na $\textbf{Własność 4., Własność 5., Własność 6.}$ widzimy, że równoważność $A=C\wedge B=D\Leftrightarrow A\times B=C\times D$ zachodzi wyłącznie dla niepustych zbiorów $A,B,C, D$ albo dla $A,B,C,D=\emptyset$. Implikacja $A=\\=C\wedge B=D\Rightarrow A\times B=C\times D$ zachodzi w ogólności(dla dowolnych zbiorów), zaś implikacja przeciwna, czyli $A\times B=C\times D\Rightarrow A=C\wedge B=D$ zachodzi wyłącznie dla niepustych zbiorów bądź gdy wszystkie zbiory $A,B,C$ i $D$ są zbiorami pustymi.
\end{remark}
\begin{property}
Niech $A$ i $B$ będą niepustymi zbiorami. W ogólności nie zachodzi: 
$A\times B=B\times A$
\end{property}
\begin{proof}[\textbf{Dowód}]
Niech $A=\{x_{1}\}$ oraz $B=\{x_{2},x_{3}\}$. Przy czym rozważane elementy są różne parami. Wówczas $A\times B=\{(x_{1},x_{2}),(x_{1},x_{3})\}$ zaś $B\times A=\{(x_{2},x_{1})$, $(x_{3},x_{1})\}$. Ponieważ $(x_{1},x_{2})\in A\times B$ i $(x_{1},x_{2})\notin B\times A$ to $A\times B\ne B\times A$. Stąd zadana równość w ogólności nie zachodzi. Zatem w ogólności iloczyn kartezjański zbiorów nie jest przemienny.
\end{proof}
\begin{property}
Prawostronna rozdzielność iloczynu kartezjańskiego względem sumy.
Dla dowolnych zbiorów $X_{1},X_{2}$ i $Y$ zachodzi $(X_{1}\cup X_{2})\times Y=\\=(X_{1}\times Y)\cup(X_{2}\times Y)$.
\end{property}
\begin{proof}[\textbf{Dowód}]
Ustalmy dowolne elementy $x$ i $y$ oraz dowolne zbiory $X_{1},X_{2}$ i $Y$. Niech $(x,y)\in(X_{1}\times Y)\cup(X_{2}\times Y)$. Wówczas z definicji wynika, że $(x,y)\in X_{1}\times Y$ lub $(x,y)\in X_{2}\times Y$. Zatem $(x\in X_{1}$ i $y\in Y)$ lub $(x\in X_{2}$ i $y\in Y)$. Stąd z rozdzielności koniunkcji względem alternatywy otrzymujemy, że $(x\in X_{1}$ lub $x\in X_{2})$ i $y\in Y$. Wówczas $x\in X_{1}\cup X_{2}$ i $y\in Y$ co z definicji daje nam $(x,y)\in(X_{1}\cup X_{2})\times Y$. Ponieważ $x$ i $y$ zostały wybrane w dowolny sposób to zachodzi inkluzja $(X_{1}\times Y)\cup(X_{2}\times Y)\subseteq(X_{1}\cup X_{2})\times Y$. $\\$
Ustalmy dowolne elementy $x$ i $y$ oraz dowolne zbiory $X_{1},X_{2}$ i $Y$. Niech $(x,y)\in(X_{1}\cup X_{2})\times Y$. Wówczas z definicji wynika, że $x\in X_{1}\cup X_{2}$ i $y\in Y$. Z definicji sumy zbiorów mamy $x\in X_{1}$ lub $x\in X_{2}$. Stąd otrzymujemy $(x\in X_{1}$ lub $x\in X_{2})$ i $y\in Y$. Zatem z rozdzielności koniunkcji względem alternatywy dostajemy, że $(x\in X_{1}$ i $y\in Y)$ lub $(x\in X_{2}$ i $y\in Y)$. Wówczas z definicji iloczynu kartezjańskiego zbiorów wynika $(x,y)\in X_{1}\times Y$ lub $(x,y)\in X_{2}\times Y$. Stąd z definicji sumy zbiorów jest $(x,y)\in(X_{1}\times Y)\cup(X_{2}\times Y)$. Ponieważ $x$ i $y$ zostały wybrane jako dowolne to zachodzi inkluzja $(X_{1}\cup X_{2})\times Y\subseteq(X_{1}\times Y)\cup(X_{2}\times Y)$. $\\$
Ponieważ zachodzą obie inkluzje to zadana równość jest w ogólności prawdziwa.
\end{proof}
\begin{property}
Lewostronna rozdzielność iloczynu kartezjańskiego względem sumy.
Dla dowolnych zbiorów $X,Y_{1},Y_{2}$ zachodzi $X\times\left (Y_{1}\cup Y_{2} \right)=\\=\left (X\times Y_{1} \right)\cup \left (X\times Y_{2} \right)$
\end{property}
\begin{proof}[\textbf{Dowód}]
Ustalmy dowolne elementy $x,y$ oraz dowolne zbiory $X,Y_{1},Y_{2}$. Niech $\left (x,y \right)\in \left (X\times Y_{1} \right)\cup \left (X\times Y_{2} \right)$. Wówczas z definicji sumy zbiorów wynika, że $\left (x,y \right)\in X\times Y_{1}$ lub $\left (x,y \right)\in X\times Y_{2}$. Następnie na mocy definicji iloczynu kartezjańskiego zbiorów dostajemy $(x\in X$ i $y\in Y_{1})$ lub $(x\in X$ i $y\in Y_{2})$. Zatem z rozdzielności koniunkcji względem alternatywy mamy $x\in X$ i $(y\in Y_{1}$ lub $y\in Y_{2})$. Wówczas z definicji sumy zbiorów jest $y\in Y_{1}\cup Y_{2}$. Ale $x\in X$ czyli z definicji iloczynu kartezjańskiego zbiorów wynika, że $\left (x,y \right)\in X\times \left (Y_{1}\cup Y_{2} \right)$. Wobec dowolności wyboru elementów $x$ i $y$ otrzymujemy $\left (x,y \right)\in \left (X\times Y_{1} \right)\cup \left (X\times Y_{2} \right)\Rightarrow \left (x,y \right)\in X\times \left ( Y_{1}\cup Y_{2}\right)$. Co jest równoważne temu, iż $\left (X\times Y_{1} \right)\cup \left (X\times Y_{2} \right)\subseteq X\times \left (Y_{1}\cup Y_{2} \right)$. $\\$
Ustalmy dowolne elementy $x,y$ oraz dowolne zbiory $X,Y_{1},Y_{2}$. Niech $\left (x,y \right)\in X\times \left (Y_{1}\cup Y_{2} \right)$. Wówczas z definicji iloczynu kartezjańskiego zbiorów wynika, że $x\in X$ i $y\in Y_{1}\cup Y_{2}$. Następnie z definicji sumy zbiorów jest $y\in Y_{1}$ lub $y\in Y_{2}$. Zatem z rozdzielności koniunkcji względem alternatywy mamy $(x\in X$ i $y\in Y_{1})$ lub $(x\in X$ i $y\in Y_{2})$ co z definicji iloczynu kartezjańskiego zbiorów daje nam, że $\left (x,y \right)\in X\times Y_{1}$ lub $\left (x,y \right)\in X\times Y_{2}$. Wówczas z definicji sumy zbiorów wynika $\left (x,y \right)\in \left (X\times Y_{1} \right)\cup \left (X\times Y_{2} \right)$. Wobec dowolności wyboru elementów $x$ i $y$ otrzymujemy $\left (x,y \right)\in X\times \left (Y_{1}\cup Y_{2} \right)\Rightarrow \left (x,y \right)\in \left (X\times Y_{1} \right)\cup \left (X\times Y_{2} \right)$. Co jest równoważne temu, iż $X\times \left (Y_{1}\cup Y_{2} \right)\subseteq \left (X\times Y_{1} \right)\cup \left (X\times Y_{2} \right)$. $\\$
Ponieważ zachodzą obie inkluzje to zadana równość jest w ogólności prawdziwa.
\end{proof}
\begin{property}
Prawostronna rozdzielność iloczynu kartezjańskiego względem różnicy.
Dla dowolnych zbiorów $X_{1},X_{2}$ i $Y$ zachodzi $\left (X_{1}\setminus X_{2} \right)\times Y=\left (X_{1}\times Y \right)\setminus \left (X_{2}\times Y \right)$.
\end{property}
\begin{proof}[\textbf{Dowód}]
Ustalmy dowolne elementy $x,y$ oraz dowolne zbiory $X_{1},X_{2}$ i $Y$. Niech $\left (x,y \right)\in \left (X_{1}\times Y \right)\setminus \left (X_{2}\times Y \right)$. Wówczas z definicji różnicy zbiorów wynika, że $\left (x,y \right)\in X_{1}\times Y$ i $\left (x,y \right)\notin X_{2}\times Y$. Następnie z definicji iloczynu kartezjańskiego zbiorów jest $x\in X_{1}$ i $y\in Y$. Zatem skoro $\left (x,y \right)\notin X_{2}\times Y$ to $x\notin X_{2}$ bo $y\in Y$. Czyli mamy $x\in X_{1}$ i $x\notin X_{2}$ i $y\in Y$ z czego wynika, że $x\in X_{1}\setminus X_{2}$ i $y\in Y$. Wówczas z definicji iloczynu kartezjańskiego zbiorów otrzymujemy $\left (x,y \right)\in \left (X_{1}\setminus X_{2} \right)\times Y$. Wobec dowolności wyboru elementów $x$ i $y$ otrzymujemy $\left (x,y \right)\in \left (X_{1}\times Y \right)\setminus \left (X_{2}\times Y \right)\Rightarrow \left (x,y \right)\in \left (X_{1}\setminus X_{2} \right)\times Y$. A jest to równoważne temu, iż $\left (X_{1}\times Y \right)\setminus \left (X_{2}\times Y \right)\subseteq \left (X_{1}\setminus X_{2} \right)\times Y$. $\\$
Ustalmy dowolne elementy $x,y$ oraz dowolne zbiory $X_{1},X_{2}$ i $Y$. Niech $\left (x,y \right)\in \left (X_{1}\setminus X_{2} \right)\times Y$. Wówczas z definicji iloczynu kartezjańskiego zbiorów wynika, że $x\in X_{1}\setminus X_{2}$ i $y\in Y$. Stąd z definicji różnicy zbiorów jest $x\in X_{1}$ i $x\notin X_{2}$. Zatem skoro $y\in Y$ to z przemienności i łączności koniunkcji mamy $(x\in X_{1}$ i $y\in Y)$ i $x\notin X_{2}$ czyli prawdą jest, że $(x\in X_{1}$ i $y\in Y)$ i $(x\notin X_{2}$ i $y\in Y)$. Wówczas z definicji iloczynu kartezjańskiego zbiorów jest $\left (x,y \right)\in X_{1}\times Y$ i $\left (x,y \right)\notin X_{2}\times Y$. Zatem z definicji różnicy zbiorów dostajemy $\left (x,y \right)\in \left (X_{1}\times Y \right)\setminus \left (X_{2}\times Y \right)$. Wobec dowolności wyboru elementów $x$ i $y$ otrzymujemy $\left (x,y \right)\in \left (X_{1}\setminus X_{2} \right)\times Y\Rightarrow \left (x,y \right)\in \left (X_{1}\times Y \right)\setminus \left (X_{2}\times Y \right)$. Jest to równoważne temu, że $\left (X_{1}\setminus X_{2} \right)\times Y\subseteq \left (X_{1}\times Y \right)\setminus \left (X_{2}\times Y \right)$. $\\$
Ponieważ zachodzą obie inkluzje to zadana równość jest w ogólności prawdziwa.
\end{proof}
\begin{property}
Lewostronna rozdzielność iloczynu kartezjańskiego względem różnicy.
Dla dowolnych zbiorów $X,Y_{1}$ i $Y_{2}$ zachodzi $X\times(Y_{1}\setminus Y_{2})=\\=(X\times Y_{1})\setminus(X\times Y_{2})$.
\end{property}
\begin{proof}[\textbf{Dowód}]
Ustalmy dowolne elementy $x,y$ oraz dowolne zbiory $X,Y_{1}$ i $Y_{2}$. Niech $(x,y)\in (X\times Y_{1})\setminus(X\times Y_{2})$ Wówczas z definicji różnicy zbiorów wynika, że $(x,y)\in X\times Y_{1}$ i $(x,y)\notin X\times Y_{2}$. Zatem z definicji iloczynu kartezjańskiego zbiorów mamy $(x\in X$ i $y\in Y_{1})$ i $(x\in X$ i $y\notin Y_{2})$. Nie może być $x\notin X$ lub $y\notin Y_{2}$ bo $x\in X$. Zatem z tego powodu jest  $x\in X$ i $y\notin Y_{2}$. Wówczas z przemienności i łączności koniunkcji dostajemy $x\in X$ i $(y\in Y_{1}$ i $y\notin Y_{2})$. Stąd $x\in X$ i $y\in Y_{1}\setminus Y_{2}$. Wówczas z definicji $(x,y)\in X\times(Y_{1}\setminus Y_{2})$. Ponieważ elementy $x$ i $y$ zostały wybrane dowolnie to zachodzi inkluzja $(X\times Y_{1})\setminus(X\times Y_{2})\subseteq X\times(Y_{1}\setminus Y_{2})$. $\\$
Ustalmy dowolne elementy $x,y$ oraz dowolne zbiory $X,Y_{1}$ i $Y_{2}$. Niech $(x,y)\in X\times(Y_{1}\setminus Y_{2})$. Wówczas z definicji $x\in X$ i $y\in Y_{1}\setminus Y_{2}$. Zatem $y\in Y_{1}$ i $y\notin Y_{2}$. Czyli otrzymujemy $(x\in X$ i $y\in Y_{1})$ i $y\notin Y_{2}$. Stąd z łączności koniunkcji $(x\in X$ i $y\in Y_{1})$ i $(x\in X$ i $y\notin Y_{2})$. Wówczas z definicji $(x,y)\in X\times Y_{1}$ i $(x,y)\notin X\times Y_{2}$ bo $y\notin Y_{2}$ a $x\in X$. Zatem $(x,y)\in(X\times Y_{1})\setminus(X\times Y_{2})$. Ponieważ $x$ i $y$ zostały wybrane w sposób dowolny to $X\times(Y_{1}\setminus Y_{2})\subseteq(X\times Y_{1})\setminus(X\times Y_{2})$. $\\$
Skoro zachodzą obie inkluzje to zadana równość jest w ogólności prawdziwa.
\end{proof}
\begin{property}
Prawostronna rozdzielność iloczynu kartezjańskiego względem przekroju.
Niech $X_{1},X_{2}$ i $Y$ będą dowolnymi zbiorami. Wówczas zachodzi: $(X_{1}\cap X_{2})\times Y=(X_{1}\times Y)\cap(X_{2}\times Y)$.
\end{property}
\begin{proof}[\textbf{Dowód}]
Ustalmy dowolne elementy $x$ i $y$ oraz dowolne zbiory $X_{1},X_{2}$ i $Y$. Niech $(x,y)\in(X_{1}\times Y)\cap(X_{2}\times Y)$. Wówczas z definicji przekroju zbiorów $(x,y)\in X_{1}\times Y$ i $(x,y)\in X_{2}\times Y$. Stąd z definicji iloczynu kartezjańskiego zbiorów $(x\in X_{1}$ i $y\in Y)$ i $(x\in X_{2}$ i $y\in Y)$. Korzystając z przemienności i łączności koniunkcji otrzymujemy $(x\in X_{1}$ i $x\in X_{2})$ i $y\in Y$. Zatem z definicji iloczynu(przekroju) zbiorów dostajemy $(x\in X_{1}\cap X_{2})$ i $y\in Y$, zaś z definicji iloczynu kartezjańskiego zbiorów $(x,y)\in(X_{1}\cap X_{2})\times Y$. Ponieważ $x$ i $y$ zostały wybrane jako dowolne elementy to zachodzi wówczas inkluzja $(X_{1}\times Y)\cap(X_{2}\times Y)\subseteq(X_{1}\cap X_{2})\times Y$. $\\$
Ustalmy dowolne elementy $x$ i $y$ oraz dowolne zbiory $X_{1},X_{2}$ i $Y$. Niech $(x,y)\in(X_{1}\cap X_{2})\times Y$. Wówczas z definicji iloczynu kartezjańskiego zbiorów wynika, że
$x\in X_{1}\cap X_{2}$ i $y\in Y$. Następnie z definicji iloczynu(przekroju) zbiorów mamy $x\in X_{1}$ i $x\in X_{2}$. Zatem otrzymujemy $(x\in X_{1}$ i $x\in X_{2})$ i $y\in Y$. Korzystając z przemienności i łączności koniunkcji dostajemy $(x\in X_{1}$ i $y\in Y)$ i $(x\in X_{2}$ i $y\in Y)$. Wówczas z definicji iloczynu kartezjańskiego zbiorów wynika, że $(x,y)\in X_{1}\times Y$ i $(x,y)\in X_{2}\times Y$, a z definicji przekroju zbiorów $(x,y)\in(X_{1}\times Y)\cap(X_{2}\times Y)$. Ponieważ $x$ i $y$ zostały wybrane w sposób dowolny to zachodzi $(X_{1}\cap X_{2})\times Y\subseteq(X_{1}\times Y)\cap(X_{2}\times Y)$. $\\$
Wobec tego, że zachodzą obie inkluzje to w ogólności zadana równość jest prawdziwa.
\end{proof}
\begin{property}
Lewostronna rozdzielność iloczynu kartezjańskiego względem przekroju.
Niech $X,Y_{1}$ i $Y_{2}$ będą dowolnymi zbiorami. Wówczas zachodzi: $X\times(Y_{1}\cap Y_{2})=(X\times Y_{1})\cap(X\times Y_{2})$.
\end{property}
\begin{proof}[\textbf{Dowód}]
Ustalmy dowolne elementy $x$ i $y$ oraz dowolne zbiory $X,Y_{1}$ i $Y_{2}$. Niech $(x,y)\in(X\times Y_{1})\cap(X\times Y_{2})$. Wówczas z definicji przekroju zbiorów wynika, że $(x,y)\in X\times Y_{1}$ i $(x,y)\in X\times Y_{2}$. Z definicji iloczynu kartezjańskiego zbiorów otrzymujemy zatem, że $(x\in X$ i $y\in Y_{1})$ i $(x\in X$ i $y\in Y_{2})$. Z łączności koniunkcji dostajemy $x\in X$ i $(y\in Y_{1}$ i $y\in Y_{2})$. Stąd $x\in X$ i $y\in Y_{1}\cap Y_{2}$. Zatem z definicji jest $(x,y)\in X\times(Y_{1}\cap Y_{2})$. Ponieważ $x$ i $y$ zostały wybrane jako dowolne elementy to zachodzi inkluzja $(X\times Y_{1})\cap(X\times Y_{2})\subseteq X\times(Y_{1}\cap Y_{2})$. $\\$
Ustalmy dowolne elementy $x$ i $y$ oraz dowolne zbiory $X,Y_{1}$ i $Y_{2}$. Niech $(x,y)\in X\times(Y_{1}\cap Y_{2})$. Wówczas z definicji wynika, że $x\in X$ i $y\in Y_{1}\cap Y_{2}$. Z definicji przekroju zbiorów dostajemy $y\in Y_{1}$ i $y\in Y_{2}$. Zatem $x\in X$ i $(y\in Y_{1}$ i $y\in Y_{2})$ co z łączności koniunkcji daje nam $(x\in X$ i $y\in Y_{1})$ i $(x\in X$ i $y\in Y_{2})$.
Stąd z definicji iloczynu kartezjańskiego zbiorów $(x,y)\in X\times Y_{1}$ i $(x,y)\in X\times Y_{2}$. Co z definicji przekroju zbiorów daje nam $(x,y)\in(X\times Y_{1})\cap(X\times Y_{2})$. Ponieważ elementy $x$ i $y$ zostały wybrane jako dowolne to zachodzi inkluzja $X\times(Y_{1}\cap Y_{2})\subseteq(X\times Y_{1})\cap(X\times Y_{2})$. Ponieważ zachodzą obie inkluzje to zadana równość jest w ogólności prawdziwa.
\end{proof}
\begin{property}
Zgodność iloczynu kartezjańskiego względem przekroju.
Niech $X_{1},X_{2},Y_{1},Y_{2}$ będą dowolnymi zbiorami. Wówczas zachodzi: $\\$ $\left (X_{1}\cap Y_{1} \right)\times \left (X_{2}\cap Y_{2} \right)=\left (X_{1}\times X_{2} \right)\cap \left (Y_{1}\times Y_{2} \right)$.
\end{property}
\begin{proof}[\textbf{Dowód}]
Ustalmy dowolne elementy $x,y$ oraz dowolne zbiory $X_{1},X_{2},Y_{1},Y_{2}$. Niech $\left (x,y \right)\in \left (X_{1}\cap Y_{1} \right)\times \left (X_{2}\cap Y_{2} \right)$. Wówczas z definicji iloczynu kartezjańskiego zbiorów wynika, że $x\in X_{1}\cap Y_{1}$ i $y\in X_{2}\cap Y_{2}$. Zatem z definicji iloczynu zbiorów jest $x\in X_{1}$ i $x\in Y_{1}$ i $y\in X_{2}$ i $y\in Y_{2}$. Następnie z przemienności i łączności koniunkcji mamy $(x\in X_{1}$ i $y\in X_{2})$ i $(x\in Y_{1}$ i $y\in Y_{2})$. Stąd z definicji iloczynu kartezjańskiego zbiorów dostajemy $\left (x,y \right)\in X_{1}\times X_{2}$ i $\left (x,y \right)\in Y_{1}\times Y_{2}$ co z definicji iloczynu zbiorów daje nam $\left (x,y \right)\in \left (X_{1}\times X_{2} \right)\cap \left (Y_{1}\times Y_{2} \right)$. Wobec dowolności wyboru elementów $x$ i $y$ otrzymujemy, że $\left (x,y \right)\in \left (X_{1}\cap Y_{1} \right)\times \left (X_{2}\cap Y_{2} \right)\Rightarrow \left (x,y \right)\in \left (X_{1}\times X_{2} \right)\cap \left (Y_{1}\times Y_{2} \right)$. Jest to równoważne temu, iż $\left [\left (X_{1}\cap Y_{1} \right)\times \left (X_{2}\cap Y_{2} \right) \right]\subseteq \left [\left (X_{1}\times X_{2} \right)\cap \left (Y_{1}\times Y_{2} \right) \right]$. $\\$ $\\$
Ustalmy dowolne elementy $x,y$ oraz dowolne zbiory $X_{1},X_{2},Y_{1},Y_{2}$. Niech $\left (x,y \right)\in \left (X_{1}\times X_{2} \right)\cap \left (Y_{1}\times Y_{2} \right)$. Wówczas z definicji iloczynu zbiorów wynika, że $\left (x,y \right)\in X_{1}\times X_{2}$ i $\left (x,y \right)\in Y_{1}\times Y_{2}$. Zatem z definicji iloczynu kartezjańskiego zbiorów jest $x\in X_{1}$ i $y\in X_{2}$ i $x\in Y_{1}$ i $y\in Y_{2}$. Następnie z przemienności i łączności koniunkcji dostajemy $(x\in X_{1}$ i $x\in Y_{1})$ i $(y\in X_{2}$ i $y\in Y_{2})$. Stąd z definicji iloczynu zbiorów mamy $x\in X_{1}\cap Y_{1}$ i $y\in X_{2}\cap Y_{2}$ co z definicji iloczynu kartezjańskiego zbiorów daje nam $\left (x,y \right)\in \left (X_{1}\cap Y_{1} \right)\times \left (X_{2}\cap Y_{2} \right)$. Wobec dowolności wyboru elementów $x$ i $y$ otrzymujemy, że $\left (x,y \right)\in \left (X_{1}\times X_{2} \right)\cap \left (Y_{1}\times Y_{2} \right)\Rightarrow \left (x,y \right)\in \left (X_{1}\cap Y_{1} \right)\times \left (X_{2}\cap Y_{2} \right)$ co jest równoważne temu, iż $\left [\left (X_{1}\times X_{2} \right)\cap \left (Y_{1}\times Y_{2} \right) \right]\subseteq \left [\left (X_{1}\cap Y_{1} \right)\times \left (X_{2}\cap Y_{2} \right) \right]$.$\\$
Ponieważ zachodzą obie inkluzje to zadana równość jest w ogólności prawdziwa.
\end{proof}
\begin{property}
Dla dowolnych zbiorów $X_{1},X_{2},Y_{1}$ i $Y_{2}$ zachodzi inkluzja $(X_{1}\times X_{2})\cup(Y_{1}\times Y_{2})\subseteq(X_{1}\cup Y_{1})\times(X_{2}\cup Y_{2})$.
\end{property}
\begin{proof}[\textbf{Dowód}]
Ustalmy dowolne zbiory $X_{1},X_{2},Y_{1}$ i $Y_{2}$ oraz dowolne elementy $x$ i $y$. Niech $(x,y)\in(X_{1}\times X_{2})\cup(Y_{1}\times Y_{2})$. Wówczas z definicji sumy zbiorów wynika, że $(x,y)\in(X_{1}\times X_{2})$ lub $(x,y)\in(Y_{1}\times Y_{2})$. Stąd jeżeli $(x,y)\in X_{1}\times X_{2}$ to z definicji iloczynu kartezjańskiego zbiorów $x\in X_{1}$ i $y\in X_{2}$. Zatem skoro $x\in X_{1}$ to z definicji sumy zbiorów $x\in X_{1}\cup Y_{1}$ nawet jeśli $x\notin Y_{1}$. Podobnie skoro $y\in X_{2}$ to z definicji sumy zbiorów $y\in X_{2}\cup Y_{2}$. Wówczas otrzymujemy $x\in X_{1}\cup Y_{1}$ i $y\in X_{2}\cup Y_{2}$ co z definicji iloczynu kartezjańskiego zbiorów daje nam $(x,y)\in(X_{1}\cup Y_{1})\times(X_{2}\cup Y_{2})$. Stąd w tym przypadku zadana inkluzja zachodzi. Analogicznie, gdy $(x,y)\in Y_{1}\times Y_{2}$. Z definicji iloczynu kartezjańskiego zbiorów wynika, że $x\in Y_{1}$ i $y\in Y_{2}$. Wówczas z definicji sumy zbiorów mamy, że $x\in X_{1}\cup Y_{1}$ oraz $y\in X_{2}\cup Y_{2}$ co z definicji iloczynu kartezjańskiego zbiorów daje nam $(x,y)\in(X_{1}\cup Y_{1})\times(X_{2}\cup Y_{2})$. Zatem i w tym wypadku zadana inkluzja zachodzi. Wówczas bez względu na to, do którego zbioru $X_{1}\times X_{2}$ czy $Y_{1}\times Y_{2}$ należy para elementów $(x,y)$ to zadana inkluzja jest prawdziwa. Stąd wobec dowolności wyboru elementów $x$ i $y$ zadana inkluzja zachodzi w ogólności.
\end{proof}
\begin{property}
Niech $X_{1},X_{2},Y_{1}$ i $Y_{2}$ będą zbiorami. W ogólności nie zachodzi równość $(X_{1}\times X_{2})\cup(Y_{1}\times Y_{2})=(X_{1}\cup Y_{1})\times(X_{2}\cup Y_{2})$.
\end{property}
\begin{proof}[\textbf{Dowód}]
Niech $X_{1}=\{x_{1},x_{2}\}$, $X_{2}=\{x_{4},x_{5}\}, Y_{1}=\{y_{1},y_{2}\}, Y_{2}=\{x_{1},y_{5}\}$, przy czym rozważane elementy są parami różne. Wówczas $(X_{1}\times X_{2})\cup(Y_{1}\times Y_{2})=[\{x_{1},x_{2}\}\times\{x_{4},x_{5}\}]\cup[\{y_{1},y_{2}\}\times\{x_{1},y_{5}\}]$ $= \{(x_{1},x_{4}),(x_{1},x_{5})$, $(x_{2},x_{4}),(x_{2},x_{5})\}$ $\cup\{(y_{1},x_{1}),(y_{1},y_{5}),(y_{2},x_{1}),(y_{2},y_{5})\}$ co daje nam zbiór $\{(x_{1}$, $x_{4}),(x_{1},x_{5}),(x_{2},x_{4})$, $(x_{2},x_{5}),(y_{1},x_{1}),(y_{1},y_{5}),(y_{2},x_{1}),(y_{2},y_{5})\}$.
Natomiast $\\$ $(X_{1}\cup Y_{1})\times(X_{2}\cup Y_{2})=\{x_{1},x_{2},y_{1},y_{2}\}\times\{x_{1},x_{4},x_{5},y_{5}\}$. Łatwo zauważyć, że para $(x_{1},x_{1})\notin(X_{1}\times X_{2})\cup(Y_{1}\times Y_{2})$ i to już wystarczy aby zadana równość nie zachodziła. To dowodzi, że iloczyn kartezjański zbiorów nie jest w ogólności rozdzielny względem ich sum.
\end{proof}
\begin{exercise}
Niech $X,Y$, $A\subseteq X$, $B\subseteq Y$ będą zbiorami. Wybierzmy dowolne elementy $x,y$. Czym jest suma $[(X\setminus A)\times Y]\cup[X\times(Y\setminus B)]$ ?
Z prawostronnej rozdzielności iloczynu kartezjańskiego względem różnicy zbiorów wynika, że $(X\setminus A)\times Y=(X\times Y)\setminus(A\times Y)$ zaś z lewostronnej rozdzielności iloczynu kartezjańskiego względem różnicy zbiorów $X\times(Y\setminus B)=(X\times Y)\setminus(X\times B)$.
Wówczas $[(X\setminus A)\times Y]\cup[X\times(Y\setminus B)]=[(X\times Y)\setminus(A\times Y)]\cup[(X\times Y)\setminus(X\times B)]$. Ale z prawa de Morgana dla różnicy zbiorów otrzymujemy, że $[(X\times Y)\setminus(A\times Y)]\cup[(X\times Y)\setminus(X\times B)]=(X\times Y)\setminus[(A\times Y)\cap(X\times B)]$. Czyli $(x,y)\in(X\times Y)\setminus[(A\times Y)\cap(X\times B)]\Leftrightarrow(x,y)\in X\times Y$ i $(x,y)\notin(A\times Y)\cap(X\times B)$. Wówczas $(x,y)\notin A\times Y$ lub $(x,y)\notin X\times B$. Zatem $(x\in X$ i $y\in Y)$ i $[(x,y)\notin A\times Y$ lub $(x,y)\notin X\times B]$. Stąd $(x\in X$ i $y\in Y)$ i $[(x\notin A$ lub $y\notin Y)$ lub $(x\notin X$ lub $y\notin B)]$. Z rozdzielności koniunkcji względem alternatywy jest
$[(x\in X$ i $y\in Y)$ i $(x\notin A$ lub $y\notin Y)]$ lub $[(x\in X$ i $y\in Y)$ i $(x\notin X$ lub $y\notin B)]$. Stąd $[(x\in X$ i $y\in Y)$ i $x\notin A)]$ lub $[(x\in X$ i $y\in Y)$ i $y\notin Y)]$ lub $[(x\in X$ i $y\in Y)$ i $(x\notin X)]$ lub $[(x\in X$ i $y\in Y)$ i $(y\notin B)]$. Z przemienności lub łączności koniunkcji $[(x\in X$ i $x\notin A)$ i $y\in Y]$ lub $[x\in X$ i $(y\in Y$ i $y\notin Y)]$ lub $[(x\in X$ i $x\notin X)$ i $y\in Y]$ lub $[x\in X$ i $(y\in Y$ i $y\notin B)]$. Pozbywamy się dwóch sprzecznych koniunkcji i otrzymujemy z definicji różnicy zbiorów $(x\in X\setminus A$ i $y\in Y)$ lub $(x\in X$ i $y\in Y\setminus B)$. Co z definicji dopełnienia i iloczynu kartezjańskiego zbiorów daje nam $(x,y)\in A^c\times Y$ lub $(x,y)\in X\times B^c$. Wówczas z definicji sumy zbiorów mamy $(x,y)\in(A^c\times Y)\cup(X\times B^c)$.
$\\$ Cofnijmy się do momentu $(x,y)\in(X\times Y)\setminus[(A\times Y)\cap(X\times B)]$. Korzystając ze zgodności pomiędzy iloczynem kartezjańskim a przekrojem zbiorów otrzymamy $(x,y)\in X\times Y\setminus(A\cap X)\times(Y\cap B)$, bo $(A\times Y)\cap(X\times B)=\\=(A\cap X)\times(Y\cap B)$, a wówczas $(x,y)\in X\times Y\setminus(A\times B)$. Bowiem z założenia mieliśmy $A\subseteq X$ i $B\subseteq Y$. Poniżej dowód tej własności pomiędzy iloczynem kartezjańskim zbiorów a dopełnieniem.
\end{exercise}
\begin{property}
Niech $X,Y$, $A\subseteq X$, $B\subseteq Y$ będą zbiorami. Wówczas zachodzi:
$(A\times B)^c=(A^c\times Y)\cup(X\times B^c)$.
\end{property}
\begin{proof}[\textbf{Dowód}]
$(A\times B)^c=(A^c\times Y)\cup(X\times B^c)$.
Ustalmy dowolne zbiory $X,Y, A\subseteq X$ i $B\subseteq Y$ oraz dowolne elementy $x$ i $y$. Niech $(x,y)\in(A^c\times Y)\cup(X\times B^c)$. Wówczas z definicji sumy zbiorów wynika, że $(x,y)\in A^c\times Y$ lub $(x,y)\in X\times B^c$. Zatem z definicji iloczynu kartezjańskiego zbiorów mamy dwa przypadki: $(x\in A^c$ i $y\in Y)$ lub $(x\in X$ i $y\in B^c)$. Jeżeli $x\in A^c$ i $y\in Y$ to $x\in X\setminus A$ i $y\in Y$, bowiem z założenia mamy $A\subseteq X$ czyli z definicji dopełnienia zbioru do przestrzeni jest $X\setminus A=A^c$. Wówczas z definicji różnicy zbiorów dostajemy $x\in X$ i $x\notin A$. Stąd z definicji iloczynu kartezjańskiego zbiorów wynika, że $(x,y)\notin A\times B$. Ale $(x,y)\in X\times Y$ bo $x\in X$ i $y\in Y$. Zatem z definicji różnicy zbiorów mamy $(x,y)\in(X\times Y)\setminus(A\times B)$. Ale $(X\times Y)\setminus(A\times B)=(A\times B)^c$ bowiem skoro $A\subseteq X$ i $B\subseteq Y$ to $A\times B\subseteq X\times Y$, a wówczas z definicji dopełnienia zbioru do przestrzeni dostajemy $(A\times B)^c=(X\times Y)\setminus(A\times B)$. Zatem $(x,y)\in(A\times B)^c$. $\\$
Analogicznie jest w drugim przypadku. Jeżeli $(x,y)\in X\times B^c$ to z definicji iloczynu kartezjańskiego zbiorów wynika, że $x\in X$ i $y\in B^c$. Wówczas z definicji dopełnienia zbioru do przestrzeni otrzymujemy $y\in Y\setminus B$, ponieważ z założenia mamy $B\subseteq Y$. Stąd $B^c=Y\setminus B$. Zatem z definicji róznicy zbiorów jest $y\in Y$ i $y\notin B$. Ale skoro $y\notin B$ to z definicji iloczynu kartezjańskiego zbiorów dostajemy $(x,y)\notin A\times B$. Ale $x\in X$ i $y\in Y$ stąd mamy $(x,y)\in X\times Y$ i $(x,y)\notin A\times B$ co z definicji róznicy zbiorów daje nam $(x,y)\in(X\times Y)\setminus(A\times B)$. Ale na mocy definicji dopełnienia zbioru do przestrzeni wynika, że $(X\times Y)\setminus(A\times B)=(A\times B)^c$. Zatem $(x,y)\in(A\times B)^c$.
Stąd wobec dowolności wyboru elementów $x$ i $y$ zachodzi inkluzja $(A^c\times Y)\cup(X\times B^c)\subseteq(A\times B)^c$. $\\$
Ustalmy dowolne zbiory $X,Y, A\subseteq X$ i $B\subseteq Y$ oraz dowolne elementy $x$ i $y$. Niech $(x,y)\in(A\times B)^c$. Z założeń wynika, że $A\times B\subseteq X\times Y$, stąd z definicji dopełnienia zbioru do przestrzeni mamy $(X\times Y)\setminus(A\times B)=(A\times B)^c$. Zatem $(x,y)\in(X\times Y)\setminus(A\times B)$. Wówczas z definicji różnicy zbiorów otrzymujemy $(x,y)\in X\times Y$ i $(x,y)\notin A\times B$ co z definicji iloczynu kartezjańskiego zbiorów daje nam $(x\in X$ i $y\in Y)$ i $(x\notin A$ lub $y\notin B)$. Korzystając z rozdzielności koniunkcji względem alternatywy dostajemy zatem $(x\in X$ i $y\in Y$ i $x\notin A)$ lub $(x\in X$ i $y\in Y$ i $y\notin B)$. Stąd z przemienności koniunkcji otrzymujemy $(x\in X$ i $x\notin A$ i $y\in Y)$ lub $(x\in X$ i $y\in Y$ i $y\notin B)$. Wówczas z definicji różnicy zbiorów wynika, że $(x\in X\setminus A$ i $y\in Y)$ lub $(x\in X$ i $y\in Y\setminus B)$. Ale jak już wiemy z dowodu pierwszej inkluzji $X\setminus A=A^c$ oraz $Y\setminus B=B^c$. Zatem $(x\in A^c$ i $y\in Y)$ lub $(x\in X$ i $y\in B^c)$. Zatem z definicji iloczynu kartezjańskiego zbiorów otrzymujemy, że $(x,y)\in A^c\times Y$ lub $(x,y)\in X\times B^c$ co z definicji sumy zbiorów daje nam $(x,y)\in(A^c\times Y)\cup(X\times B^c)$. A ponieważ $x$ i $y$ zostały wybrane jako dowolne elementy to inkluzja $(A\times B)^c\subseteq(A^c\times Y)\cup(X\times B^c)$ jest prawdziwa. $\\$
Wówczas skoro zachodzą obie inkluzje to zadana równość jest w ogólności prawdziwa.
\end{proof}
\begin{property}
Dla dowolnych zbiorów $X,Y,A\subseteq X$ i $B\subseteq Y$ zachodzi równość: $(A\times B)^c=(A^c\times B)\cup(A\times B^c)\cup(A^c\times B^c)$.
\end{property}
\begin{proof}[\textbf{Dowód}]
Ustalmy dowolne zbiory $X,Y, A\subseteq X$ i $B\subseteq Y$ oraz dowolne elementy $x$ i $y$. Niech $(x,y)\in(A^c\times B)\cup(A\times B^c)\cup(A^c\times B^c)$. Wówczas z definicji sumy zbiorów wynikają trzy przypadki: $(x,y)\in A^c\times B$ lub $(x,y)\in A\times B^c$ lub $(x,y)\in A^c\times B^c$. Przeanalizujmy każdy po kolei. $\\$
Jeżeli $(x,y)\in A^c\times B$ to z definicji iloczynu kartezjańskiego zbiorów otrzymujemy $x\in A^c$ i $y\in B$. Ponieważ $A\subseteq X$ to z definicji dopełnienia zbioru do przestrzeni wynika, że $A^c=X\setminus A$. Zatem $x\in X\setminus A$ co z definicji różnicy zbiorów daje nam $x\in X$ i $x\notin A$. Ponieważ $y\in B$ to otrzymujemy zatem $x\in X$ i $x\notin A$ i $y\in B$. Korzystając z łączności koniunkcji mamy $(x\in X$ i $x\notin A)$ i $y\in B$ co daje nam $(x\in X$ i $y\in B)$ i $(x\notin A$ i $y\in B)$. Ale z założenia jest $B\subseteq Y$ czyli jeżeli $y\in B$ to $y\in Y$. Stąd $(x\in X$ i $y\in Y)$ i $(x\notin A$ i $y\in B)$. Zatem z definicji iloczynu kartezjańskiego zbiorów dostajemy, że $(x,y)\in X\times Y$ i $(x,y)\notin A\times B$. Wówczas z definicji różnicy zbiorów jest $(x,y)\in(X\times Y)\setminus(A\times B)$. Ponieważ z założeń wynika, iż $A\times B\subseteq X\times Y$ to z definicji dopełnienia zbioru do przestrzeni otrzymujemy $(X\times Y)\setminus(A\times B)=(A\times B)^c$. Stąd $(x,y)\in(A\times B)^c$. Zatem w tym przypadku zachodzi inkluzja $(A^c\times B)\cup(A\times B^c)\cup(A^c\times B^c)\subseteq(A\times B)^c$. $\\$
Jeżeli $(x,y)\in A\times B^c$ to z definicji iloczynu kartezjańskiego zbiorów wynika, że $x\in A$ i $y\in B^c$. Z założenia $B\subseteq Y$ i z definicji dopełnienia zbioru do przestrzeni mamy $B^c=Y\setminus B$. Stąd $y\in Y\setminus B$ co z definicji różnicy zbiorów daje $y\in Y$ i $y\notin B$. Zatem dostajemy $x\in A$ i $y\in Y$ i $y\notin B$. Korzystając z łączności koniunkcji otrzymujemy wówczas $(x\in A$ i $y\in Y)$ i $(x\in A$ i $y\notin B)$. Ale z założenia $A\subseteq X$ wynika, że jeżeli $x\in A$ to $x\in X$. Stąd $(x\in X$ i $y\in Y)$ i $(x\in A$ i $y\notin B)$. Stąd z definicji iloczynu kartezjańskiego zbiorów jest $(x,y)\in X\times Y$ i $(x,y)\notin A\times B$ co z definicji różnicy zbiorów daje nam $(x,y)\in(X\times Y)\setminus(A\times B)$. Ale z założeń i z definicji dopełnienia zbioru do przestrzeni mamy, że zbiór $X\times Y$ jest przestrzenią dla zbioru $(X\times Y)\setminus(A\times B)$. Zatem $(X\times Y)\setminus(A\times B)=(A\times B)^c$, a wówczas $(x,y)\in(A\times B)^c$. Tak więc w tym przypadku zachodzi inkluzja $(A^c\times B)\cup(A\times B^c)\cup(A^c\times B^c)\subseteq(A\times B)^c$. $\\$
Jeżeli $(x,y)\in A^c\times B^c$ to z definicji iloczynu kartezjańskiego zbiorów wynika, że $x\in A^c$ i $y\in B^c$. Z założeń $A\subseteq X$ i $B\subseteq Y$ oraz z definicji dopełnienia zbioru do przestrzeni otrzymujemy wówczas $x\in X\setminus A$ i $y\in Y\setminus B$. Zatem z definicji różnicy zbiorów dostajemy $x\in X$ i $x\notin A$ i $y\in Y$ i $y\notin B$. Korzystając z przemienności koniunkcji mamy $(x\in X$ i $y\in Y)$ i $(x\notin A$ i $y\notin B)$ co z definicji iloczynu kartezjańskiego zbiorów daje nam $(x,y)\in X\times Y$ i $(x,y)\notin A\times B$. Stąd z definicji różnicy zbiorów wynika, że $(x,y)\in(X\times Y)\setminus(A\times B)$. Ponieważ z założeń mamy $A\times B\subseteq X\times Y$ to z definicji dopełnienia zbioru do przestrzeni jest $(X\times Y)\setminus(A\times B)=(A\times B)^c$. Wówczas $(x,y)\in(A\times B)^c$. Zatem w tym przypadku zachodzi inkluzja $(A^c\times B)\cup(A\times B^c)\cup(A^c\times B^c)\subseteq(A\times B)^c$. $\\$
Zauważmy wówczas, że otrzymujemy alternatywę inkluzji $A^c\times B\subseteq(A\times B)^c$ lub $A\times B^c\subseteq(A\times B)^c$ lub $A^c\times B^c\subseteq(A\times B)^c$, ponieważ $(x,y)\in(A^c\times B)\cup(A\times B^c)\cup(A^c\times B^c)\Rightarrow(x,y)\in(A\times B)^c$. Stąd wobec dowolności wyboru elementów $x$ i $y$ w ogólności zachodzi inkluzja $(A^c\times B)\cup(A\times B^c)\cup(A^c\times B^c)\subseteq(A\times B)^c$.$\\$
Ustalmy dowolne zbiory $X,Y, A\subseteq X$ i $B\subseteq Y$ oraz dowolne elementy $x$ i $y$. Niech $(x,y)\in(A\times B)^c$. Wówczas z założeń i definicji dopełnienia zbioru do przestrzeni wynika, że $(x,y)\in(X\times Y)\setminus(A\times B)=(A\times B)^c$. Stąd z definicji różnicy zbiorów otrzymujemy $(x,y)\in X\times Y$ i $(x,y)\notin A\times B$. Zatem z definicji iloczynu kartezjańskiego zbiorów dostajemy $(x\in X$ i $y\in Y)$ i $(x\notin A$ lub $y\notin B)$. Wówczas $(x\in X$ i $y\in Y)$ i $[(x\in A$ i $y\notin B)$ lub $(x\notin A$ i $y\in B)$ lub $(x\notin A$ i $y\notin B)]$.$\\$
Jeżeli $(x\in X$ i $y\in Y)$ i $(x\in A$ i $y\notin B)$ to z przemienności koniunkcji dostajemy $(x\in X$ i $x\in A)$ i $(y\in Y$ i $y\notin B)$. Ponieważ z założenia $A\subseteq X$ to $(x\in X$ i $x\in A)$ jest równoważne $x\in A$, co wynika z definicji przekroju zbiorów. Zatem otrzymujemy $x\in A$ i $(y\in Y$ i $y\notin B)$. Stąd z definicji różnicy zbiorów wynika, że $x\in A$ i $y\in Y\setminus B$. Ale z założenia i definicji dopełnienia zbioru do przestrzeni mamy $y\in B^c$. Wówczas dostajemy $x\in A$ i $y\in B^c$ co z definicji iloczynu kartezjańskiego zbiorów daje nam $(x,y)\in A\times B^c$. Zatem $(A\times B)^c\subseteq(A^c\times B)\cup(A\times B^c)\cup(A^c\times B^c)$. $\\$
Jeżeli $(x\in X$ i $y\in Y)$ i $(x\notin A$ i $y\in B)$ to z przemienności koniunkcji otrzymujemy, że $(x\in X$ i $x\notin A)$ i $(y\in Y$ i $y\in B)$. Ponieważ z założenia jest $B\subseteq Y$ to $(y\in Y$ i $y\in B)$ jest równoważne $y\in B$. Zatem mamy $(x\in X$ i $x\notin A)$ i $y\in B$ co z definicji różnicy zbiorów daje nam $x\in X\setminus A$ i $y\in B$. Ale z założenia i z definicji dopełnienia zbioru do przestrzeni dostajemy $x\in A^c$, bo $X\setminus A=A^c$. Stąd jest $x\in A^c$ i $y\in B$. Wówczas z definicji iloczynu kartezjańskiego zbiorów wynika, że $(x,y)\in A^c\times B$. Zatem zachodzi inkluzja $(A\times B)^c\subseteq(A^c\times B)\cup(A\times B^c)\cup(A^c\times B^c)$.$\\$
Jeżeli $(x\in X$ i $y\in Y)$ i $(x\notin A$ i $y\notin B)$ to z przemienności koniunkcji mamy $(x\in X$ i $x\notin A)$ i $(y\in Y$ i $y\notin B)$. Wówczas z definicji różnicy zbiorów wynika, że $x\in X\setminus A$ i $y\in Y\setminus B$. Stąd z założeń i definicji dopełnienia zbioru do przestrzeni otrzymujemy $x\in A^c$ i $y\in B^c$ co z definicji iloczynu kartezjańskiego zbiorów daje nam $(x,y)\in A^c\times B^c$. Zatem zachodzi inkluzja $(A\times B)^c\subseteq(A^c\times B)\cup(A\times B^c)\cup(A^c\times B^c)$. $\\$
Wówczas w każdym z przypadków zachodzi inkluzja $(A\times B)^c\subseteq(A^c\times B)\cup(A\times B^c)\cup(A^c\times B^c)$. $\\$
Wobec tego, że zachodzą obie inkluzje przeciwne to zadana równość jest w ogólności prawdziwa.
\end{proof}